\documentclass[a4paper,11pt]{article}
\usepackage{exam-style-842-unit}
\examinehead{842 信息化概论 · 第 4 单元练习册}{练习二\quad 开发架构与 Web 服务}
\begin{document}

\unitcover{842 信息化概论 · 单元练习册}{第 4 单元\quad 软件技术与开发架构}{练习二\quad 开发架构与 Web 服务（共 100 分）}{约 70 分钟}{C/S 与 B/S 架构；三层架构与中间件；Web 服务与 SOA 体系架构；SOAP、WSDL、UDDI；HTTP 协议与网页请求—响应过程。说明：教材未设 B/S、C/S 与 Web 服务专节，本练习按考试大纲“二、4”与通用规范知识补充，答案册逐题注明依据。}

\begin{knowmap}
\textbf{两种基本架构：}C/S（客户机/服务器，客户端安装专用程序，承担表示与部分业务逻辑）$\leftrightarrow$ B/S（浏览器/服务器，浏览器即客户端，业务与数据集中在服务器）\\[2pt]
\textbf{分层：}单层 $\rightarrow$ 两层 C/S $\rightarrow$ 三层（表示层—业务逻辑层—数据访问层）$\rightarrow$ 多层/微服务；层间通过接口松耦合，中间件提供通信、事务、安全等公共服务\\[2pt]
\textbf{Web 服务：}以标准 Web 协议与中立数据格式，让不同平台、不同语言的应用互相调用，是 SOA 的主要实现技术\\[2pt]
\textbf{三要素：}SOAP（基于 XML 的消息协议）＋ WSDL（服务描述）＋ UDDI（服务注册与发现）\\[2pt]
\textbf{三角色：}服务提供者 $\rightarrow$ 服务注册中心（发布/查找）$\leftarrow$ 服务请求者（绑定调用）\\[2pt]
\textbf{网页请求—响应：}输入 URL $\rightarrow$ DNS 解析 $\rightarrow$ 建立 TCP 连接 $\rightarrow$ 发送 HTTP 请求 $\rightarrow$ 服务器处理（Web 服务器／应用服务器／数据库）$\rightarrow$ 返回 HTTP 响应 $\rightarrow$ 浏览器解析渲染
\end{knowmap}

\begin{keyterms}
\textbf{C/S 架构}：客户机/服务器架构，客户端安装专用程序并承担表示与部分业务逻辑，服务器管理数据与共享资源。\\
\textbf{B/S 架构}：浏览器/服务器架构，以 Web 浏览器为统一客户端，业务逻辑与数据集中在服务器端。\\
\textbf{三层架构}：由表示层、业务逻辑层、数据访问层组成的软件分层结构，层间通过接口交互。\\
\textbf{表示层}：与用户交互的界面层，负责信息展示、输入采集与基本校验。\\
\textbf{业务逻辑层}：实现业务规则与业务流程、提供业务服务的中间层。\\
\textbf{数据访问层}：封装对数据库等数据源的访问操作，为上层提供统一数据接口。\\
\textbf{中间件}：位于操作系统与应用软件之间、提供通信、事务、安全等公共服务的软件。\\
\textbf{Web 服务}：基于标准 Web 协议、以中立格式交换消息，实现跨平台跨语言应用互操作的分布式计算技术。\\
\textbf{SOA}：面向服务的体系架构，把业务功能封装为可复用、可组合的服务，通过接口契约松耦合集成。\\
\textbf{SOAP}：简单对象访问协议，基于 XML 的消息格式，通常绑定 HTTP 传输。\\
\textbf{WSDL}：Web 服务描述语言，用 XML 描述服务的接口、操作、消息与访问地址。\\
\textbf{UDDI}：统一描述、发现与集成，提供服务注册、发布与查找的机制。\\
\textbf{HTTP}：超文本传输协议，采用请求—响应模式，是浏览器与 Web 服务器之间的应用层协议。\\
\textbf{URL}：统一资源定位符，由协议、主机（域名或 IP）、端口、路径、查询参数等部分组成。
\end{keyterms}

\begin{contrast}{C/S / B/S（每行依次为：客户端·部署与升级·典型适用场景）}
\citem{C/S}{需安装专用客户端程序，客户端承担表示与部分业务逻辑}{客户端与服务器都要升级，维护成本较高，跨平台性差}{局域网内、交互复杂、数据量大、实时性要求高的业务系统}
\citem{B/S}{只需浏览器，客户端零安装，业务逻辑集中在服务器}{升级维护集中在服务器端，部署快、易扩展}{互联网上面向大量用户的信息发布、电商、在线办事等应用}
\end{contrast}

\newpage
\begin{contrast}{两层 C/S / 三层架构（每行依次为：结构·优点·不足）}
\citem{两层 C/S}{客户端（表示与业务逻辑）＋数据库服务器（数据管理）}{结构简单、交互响应快、开发部署容易}{业务逻辑与界面耦合，客户端“胖”，升级维护困难，安全性较弱}
\citem{三层架构}{表示层＋业务逻辑层＋数据访问层}{职责分离、层间松耦合，可维护性、可复用性与可扩展性好}{层次增多，调用链变长，性能开销与设计复杂度上升}
\end{contrast}

\begin{contrast}{SOAP / WSDL / UDDI（每行依次为：作用·主要内容·通俗理解）}
\citem{SOAP}{定义服务之间如何传递消息}{基于 XML 的消息格式与处理规则，通常绑定 HTTP}{像“信封与信纸”，规定消息怎么写、怎么送}
\citem{WSDL}{描述服务能做什么、怎么调用}{用 XML 描述服务接口、操作、消息格式与访问端点}{像“产品说明书”，告诉调用者接口与参数}
\citem{UDDI}{让服务可以被发布、查找和发现}{服务注册中心的注册、查询与发现机制}{像“电话黄页”，按需查找到服务地址}
\end{contrast}

\begin{contrast}{静态网页 / 动态网页（每行依次为：生成方式·内容特点·常见技术）}
\citem{静态网页}{服务器直接返回预先写好的 HTML 文件}{内容固定，修改需改源文件，响应快}{HTML＋CSS，纯前端}
\citem{动态网页}{服务器端程序根据请求与数据实时生成页面}{内容随用户、时间、数据变化，可与数据库交互}{JSP、PHP、ASP.NET、Node.js、Python 等}
\end{contrast}

\qtypename{一、单项选择题}{10 题，每题 3 分，共 30 分}
\q 采用 Web 浏览器作为统一客户端、把系统维护升级集中在服务器端的软件架构是（\quad）
\choicesline{A. C/S 架构}{B. B/S 架构}{C. 单机架构}{D. 主机—终端架构}

\q 在 C/S 架构中，客户端通常承担（\quad）
\choicesline{A. 只负责显示}{B. 表示逻辑和部分业务逻辑}{C. 只负责数据管理}{D. 只负责网络路由}

\q 与 C/S 架构相比，B/S 架构最突出的优点是（\quad）
\choices{A. 客户端零安装、升级维护集中在服务器端}{B. 客户端必须安装专用软件}
{C. 完全不受网络带宽影响}{D. 全部业务逻辑都在浏览器中执行}

\q 下列属于三层架构中业务逻辑层职责的是（\quad）
\choicesline{A. 页面展示与输入校验}{B. 业务规则与业务流程处理}{C. 数据库连接与 SQL 执行}{D. 网络报文的路由转发}

\q Web 服务三要素中，用于描述服务接口、操作和访问地址的是（\quad）
\choicesline{A. SOAP}{B. WSDL}{C. UDDI}{D. HTTP}

\q 在 Web 服务体系中，负责“服务注册与发现”的机制是（\quad）
\choicesline{A. SOAP}{B. WSDL}{C. UDDI}{D. SSL}

\q SOAP 消息通常采用（\quad）编码，并绑定在 HTTP 等协议上传输。
\choicesline{A. 二进制}{B. XML}{C. JSON}{D. HTML}

\q 在浏览器地址栏输入网址后，浏览器首先进行的操作是（\quad）
\choices{A. 直接向 Web 服务器发送 HTTP 请求}{B. 解析 URL 并进行 DNS 域名解析}
{C. 渲染页面并执行脚本}{D. 建立数据库连接}

\q HTTP 响应报文中的状态码 404 表示（\quad）
\choicesline{A. 请求成功}{B. 请求的资源不存在}{C. 服务器内部错误}{D. 需要身份认证}

\q 下列不属于 Web 服务特点的是（\quad）
\choices{A. 跨平台、跨语言调用}{B. 基于标准协议、松耦合}
{C. 只能在同一个局域网内调用}{D. 以服务为单位封装业务功能}

\qtypename{二、多项选择题}{3 题，每题 4 分，共 12 分；每题有两个或两个以上正确项}
\q B/S 架构相比 C/S 架构的优点包括（\quad）
\choices{A. 客户端零安装、零维护}{B. 升级维护集中在服务器端}
{C. 跨平台、易扩展}{D. 适合在互联网上向大量用户提供服务}

\q Web 服务的三要素（核心技术）包括（\quad）
\choicesline{A. SOAP}{B. WSDL}{C. UDDI}{D. HTML}

\q 网页请求与响应过程中，可能发生的有（\quad）
\choices{A. 浏览器解析 URL 并进行 DNS 域名解析}{B. 浏览器与 Web 服务器建立 TCP 连接}
{C. 浏览器向服务器发送 HTTP 请求报文}{D. 服务器返回状态码与响应体}

\qtypename{三、判断题}{6 题，每题 2 分，共 12 分；正确的画“√”，错误的画“×”并写出改正后的表述}
\q B/S 架构中客户端只需浏览器，不必安装专用客户端软件，升级维护主要集中在服务器端。\quad（\quad）\quad 改正：\fixblank

\q C/S 架构的系统升级时只需升级服务器，客户端不必更新。\quad（\quad）\quad 改正：\fixblank

\q Web 服务只能在同一编程语言、同一操作系统平台的应用之间互相调用。\quad（\quad）\quad 改正：\fixblank

\q SOAP 采用 XML 描述消息，通常通过 HTTP 协议传输。\quad（\quad）\quad 改正：\fixblank

\q WSDL 负责服务的注册与发现。\quad（\quad）\quad 改正：\fixblank

\q 网页请求与响应中，服务器返回的响应内容不一定是 HTML，也可以是 JSON、XML、图片等数据。\quad（\quad）\quad 改正：\fixblank

\qtypename{四、填空题}{6 空，每空 1 分，共 6 分}
\q C/S 指客户机/服务器架构，B/S 指\fillblanks[3.4cm]架构。

\q Web 服务的三种核心技术是 SOAP、\fillblanks[1.8cm]和 UDDI。

\q 三层架构由表示层、\fillblanks[2.6cm]和数据访问层组成。

\q HTTP 请求报文由请求行、\fillblanks[1.8cm]和请求体（实体）组成。

\q HTTP 状态码中，404 表示\fillblanks[2.8cm]。

\q 网页请求与响应过程中，浏览器把域名解析为 IP 地址，这一过程称为\fillblanks[2.2cm]解析。

\qtypename{五、简答题}{4 题，每题 6 分，共 24 分}
\q 简述 C/S 与 B/S 两种架构的含义，并比较二者的主要区别与适用场景。\anslines{2.4cm}

\q 什么是 Web 服务？简述 Web 服务体系的三种角色和三项核心技术。\anslines{2.4cm}

\q 简述三层架构的组成及各层职责，并说明分层的主要好处。\anslines{2.4cm}

\q 简述从在浏览器输入网址到页面呈现的网页请求与响应过程。\anslines{2.4cm}

\newpage
\qtypename{六、综合题}{1 题，共 16 分}
\q 某农业产业化龙头企业原有一套 C/S 架构的进销存系统，只能在办公电脑上安装使用。企业计划把系统改造为 B/S 架构，并把库存查询、订单提交等能力以 Web 服务方式开放给合作社的微信小程序和第三方电商平台。请回答：
（1）说明由 C/S 改造为 B/S 的理由（至少 3 点），并指出 B/S 相比 C/S 的不足；
（2）按三层架构说明改造后系统的分层与各层职责；
（3）说明对外开放的 Web 服务需要哪三项核心技术、各自起什么作用；
（4）描述小程序调用该 Web 服务时的请求与响应过程，并指出至少 2 项安全或性能方面的注意事项。\anslines{5cm}

\end{document}
